% ─── Conclusion ───────────────────────────────────────────────────────────────
\label{sec:conclusion}

%── 7.1 Summary of Findings ────────────────────────────────────────────────────
\subsection{Summary of Key Findings}

This section answers RQ1 and RQ2 as stated in Section~\ref{sec:rq}; RQ3 is answered in
Section~\ref{sec:power}. Each figure below is derived where it is first reported.

On RQ1 the verdict runs one way across every arena, and the reason is the same in each.
Against the fossil incumbent it is settled. The best heat pump is cheaper than a gas
boiler in 24 of the 29 markets by 2050, a count that runs from 19 to all 29 across the
carbon-price paths because only the gas boiler pays that price. Against hydrogen the
base-load contest settles too, with each carrier charged its own equipment capital, its
own last-mile network and the seasonal storage its own load shape requires. The best heat
pump leads in 22 of 29, at 2050 medians of \euro{}118.5/MWh against \euro{}122.8/MWh
and a median country gap of \euro{}15.4/MWh in its favour, and hydrogen holds the
remaining seven, five of
them salt-cavern markets.

That concentration is the paper's central structural result, and its cause divides.
Heating draws about
three-quarters of its energy between October and March while electrolytic production is
far flatter, so the molecule must be carried between seasons, at about \euro{}1.5/kg
where salt caverns exist and \euro{}3.0/kg where they do not. Hydrogen's base-load case
concentrates where that store is cheap, which is the same condition that decides the
peaking arena, so the two tests turn on one physical quantity. That shared parameter is
also why their agreement is a property of the model as much as an independent
confirmation. Geology does not carry the whole of the concentration. Charging every
country the cavern rate leaves the heat pump ahead in 19 of 29 and the non-cavern rate
leaves it ahead in 24, so the geological binary is worth five markets and decides two
of the seven exceptions. The other five turn on the retail tariff the heat pump pays.

What hydrogen keeps in those seven is not a technology result. It follows from the fiscal
treatment of the two carriers. Pricing the molecule as blue hydrogen at the \euro{}81/MWh of our own blue route \citep{IEAGHG2022blue}, the provenance one dedicated assessment concludes heating hydrogen would have \citep{Dickel2024ET29}, removes the lead in all seven; equalising the fiscal treatment
would widen the heat pump's margin in the same direction, though the levy-equalised
comparison that would size the effect is one we did not run
(Appendix~\ref{sec:limitations}). The base-load count moves with the hydrogen price path and not
with the policy scenario, from 13 markets on the fastest cost decline to all 29 if
hydrogen costs stall, because neither appliance pays a carbon price.

Where hydrogen's dispatch case is strongest it fails the investment test. Repriced as a
winter peaker it wins on operating cost, under the most carbon-ambitious H2~Push
scenario, in seven markets on the even-handed efficiency series and 15 on the evolving
one, and those wins are decisive in only the seven
salt-cavern markets of the north-west cluster that carries the firm-capacity tip, the
remainder marginal (Section~\ref{sec:power:share}). In no modelled country-scenario does the
inframarginal rent recover the peaker's capital. That is an economic signal and not a
reliability one, since the system stays adequate at a modelled loss of load averaging
about two hours a year. The case against hydrogen for heat therefore rests on the
base-load and investment arenas together, and on the same endowment in each.

Three extensions test these results without altering them. An explicit
unit-commitment cross-check of the firm fleet and an endogenous Swiss integration with
electricity-price feedback both bear on the peaking verdict, and a myopic
rolling-horizon test bears on the cost-optimal retrofit share.

On RQ2 the Swiss base-load and peaking comparisons both run against hydrogen. The best
Swiss heat pump ends \euro{}5.6/MWh below the hydrogen boiler on the uncoupled arm we
quote as the conservative one, and \euro{}9.6/MWh below it on the endogenous arm.
Switzerland has short pipeline distances to the major European producers and dear retail
electricity, but no salt caverns, so its heating load would be stored at the higher rate,
and that charge holds the ordering. The peaking comparison points the same way with far
more room, since hydrogen wins no Swiss winter peak in any scenario once the Swiss
reservoir fleet is credited as firm winter capacity. Switzerland accordingly
makes best use of the European build-out by directing it, as federal strategy already
does, to industry and mobility, and meeting residential heat through electrification
supported by its hydropower flexibility.

Common to both questions is an emissions divergence of about 342~MtCO$_2$/yr between
drift and net zero that the four scenarios do not close. From a common
\mbox{$\approx$644~MtCO$_2$} 2025 baseline they reach 352 (Current Policies), 231 (Stated Policies), 123 (H2 Push) and 10 (Net Zero)~MtCO$_2$ by
2050. Because the same technologies are available in every scenario, that gap measures
ambition.

The same divergence appears on the cost side. Each pathway carries a distinct
network-infrastructure bill cumulative to 2050, \euro{}249~bn under Current Policies,
\euro{}325~bn under Stated Policies, \euro{}412~bn under H2 Push and \euro{}506~bn under
Net Zero, dominated in the deep pathway by district-heating expansion at \euro{}356~bn
against \euro{}119~bn of electricity reinforcement. The hydrogen-push variant spends
\euro{}74~bn on a dedicated hydrogen distribution network and still requires
\euro{}89~bn of reinforcement, so hydrogen investment does not displace grid expansion
and the two networks must co-develop even in the hydrogen-heavy case.

%── 7.2 Contribution ───────────────────────────────────────────────────────────
\subsection{Contribution}

The contribution of this paper is threefold.

First, it provides what is, to our
knowledge, the first merit-order assessment of hydrogen's
role in residential heating across the EU-27, the United Kingdom and Switzerland
(29 countries, 2025 to 2050), built on a bottom-up heat-demand reconstruction
resolved to NUTS3, in which the heating share of
hydrogen is bounded above by where it can win the cold-snap peak, not assumed.
Building demand is reconstructed bottom-up from EUBUCCO footprints and TABULA
intensities, validated against Hotmaps, and carried through a 200-draw Monte
Carlo over four scenarios and an independent least-cost optimisation.

Second, it prices hydrogen the way a system operator does, on short-run marginal cost in
a dispatch, and tests it in three linked arenas, the building winter peak, the
district-heating stack and the power-sector peak. Testing all three is
what makes the exercise informative, because all three converge on a single condition.
Hydrogen can be the cheaper appliance on base-load cost, and it can win the winter-peak
dispatch, and in both tests the condition that decides it is the cost of holding a
winter-weighted load in store from one season to the next. Cheap salt-cavern storage is
what makes that condition affordable, so the same endowment is decisive in each. They are not the same seven countries in each.

Five of the seven base-load holds are cavern markets, while the seven peaking wins are
exactly the cavern set, and the two lists differ in Belgium and Ireland, which hold base
load on very high retail tariffs, and in France and Romania, which have caverns yet lose
the base-load comparison. Even in its most favourable niche a hydrogen peaker cannot
recover its capital from market rent, because the scarcity hours on which it earns are
too few. That is the missing-money problem. That convergence across arenas, shown across
29 countries and four scenarios, is what distinguishes this approach from single-arena
assessments, which return whichever verdict the arena they chose happens to carry and
cannot tell whether the arenas share a mechanism.

Third, the analysis isolates a finding of wider import. In the dispatch arenas it is
the carbon price on the fossil competitor, more than the price of hydrogen, that
opens what niche hydrogen has, while the base-load arena, in which neither clean
option pays carbon at all, turns on the cost of seasonal storage first and on how
the two carriers are taxed at the meter second. On the least-cost side the ambition
is cheap without being free. The
present-value system cost moves by 0.05 per cent between a 75 and a 100 per cent
reduction cap, so decarbonising scope-1 residential heat is close to costless at the
system level, though the cap does bind in 4 of 29 countries by 2050 (Cyprus, Croatia,
Poland and the United Kingdom) at a shadow price of up to \euro{}81/tCO$_2$, and in
7 countries at some point on the path, which is where an implicit carbon price beyond
ETS2 would have to be found. That reading holds under both perfect foresight
and the myopic rolling-horizon test (Appendix~\ref{sec:costopt}).

%── 7.3 Policy Implications ─────────────────────────────────────────────────────
\subsection{Policy Implications}
\label{sec:conc:policy}

Five policy implications follow from these findings.

\textbf{1. Policy should make heat-pump and district-heat deployment the primary
near-term priority.} Against the fossil incumbent the case is already made in 2025
price data, heat pumps being cheaper than a gas boiler in 16 of 29 countries then and
in 24 by 2050, and ETS2 widens that margin. The constraints are partly economic and
partly not. Equipment capital is now the largest single fixed item in a heat pump's
cost, at about 28 per cent of its 2050 levelised cost, which is a heavier capital
burden than either the gas or the hydrogen boiler carries, and reinforcing low- and
medium-voltage networks for the winter heat-pump peak adds between \euro{}1.1 and
38.2/MWh depending on how much of that cost the load is asked to bear. Installer
capacity, access to upfront capital and regulatory inertia sit on top of those.
Policy is therefore best directed at skilled-installer training, subsidised
financing (for example through the REPowerEU home-renovation programme),
streamlined planning permission for residential heat pumps, and at bringing down the
delivered capital and network cost to which the levelised comparison is most
sensitive.

\textbf{2. Policy should allow ETS2 to pass through to end users, conditional on
empirical measurement of the household switching response, to preserve the
switching incentive.} This paper quantifies the carbon-cost differential between
gas and heat-pump heating (an adder of \euro{}11/MWh per \euro{}50/tCO$_2$ against a median \euro{}0.7/MWh by
2030) but does not model the behavioural elasticity of the consumer response to
it, a gap that future work should close; the implication that follows is therefore
a logical one conditioned on that response being non-trivial. On that condition,
because energy-bill caps blunt the differential signal, measures that shield
residential consumers from ETS2 in the aggregate, such as broad energy-bill caps
and untargeted rebates, would slow the transition. Targeted social protection for
fuel-poor households is appropriate and necessary, but it should be designed to
preserve the switching incentive, not to suppress the carbon price across
the board.

\textbf{3. Policy should require a capacity-payment design and a co-benefit
justification for hydrogen in buildings, and should not rest the decision on the
headline cost comparison in either direction.} Where hydrogen still comes out cheaper
than a heat pump on delivered annual cost, in those seven markets, the margin comes from
the levies and network charges loaded onto retail electricity and not onto the molecule,
so it is an artefact of tax policy that tax policy can remove, and it would survive
neither an equalisation of the two treatments nor a shift of heating hydrogen onto a
blue supply route. Where hydrogen looks strongest in dispatch, in the cavern-endowed
countries at a high carbon price, no private investor recovers the peaker's capital from
market rents alone.

Any deliberate hydrogen-for-heat deployment would therefore require explicit
justification on grounds beyond the headline comparison, such as co-benefits in
particular geographies or network configurations (for example regions where major
district-heating expansion is infeasible), and should be paired with a standing capacity
payment or equivalent support to signal the absence of revenue-market recovery. These
co-benefits are not quantified in the present model and are noted as qualitative
considerations and not modelled results. The corollary on the tax side is that the levy
structure on residential electricity is itself a heating-technology policy, since
reforming it would move the base-load ranking in a large group of countries. Policy
should make these distinctions explicit, not present hydrogen as a general-purpose
decarbonisation solution for buildings.

\textbf{4. Policy should give Eastern European countries tailored transition support.}
Countries with coal-heavy electricity grids and constrained fiscal capacity, such as
Poland, the Czech Republic, Romania and Bulgaria \citep{EMBER2024electricity,
EEA2023fit55}, face the most demanding transition, since heat pumps are cost-competitive
but carry a higher near-term grid carbon intensity, and boiler bans risk imposing
disproportionate adjustment costs. This is also the region where the cost evidence most
strongly favours electrification, the heat pump's 2050 margin over the hydrogen boiler
being widest across the warm south and east, at \euro{}99/MWh in Malta, \euro{}88/MWh in
Cyprus, \euro{}61/MWh in Hungary, \euro{}48/MWh in Greece and \euro{}41/MWh in Bulgaria.
Bulgaria, the Czech Republic and Hungary have no salt-cavern endowment to make a
seasonal hydrogen store cheap, so in those markets there is no base-load case for a
hydrogen alternative to hedge against, and support can be directed squarely at the grid
and the installation base.

Poland and Romania are the exceptions within the group and should be treated separately.
Both hold caverns, and Poland is one of the seven markets in which the hydrogen boiler
still leads on delivered annual cost, though by only \euro{}2/MWh, a margin narrower
than the fiscal wedge between the two carriers. Romania, despite the same geology, sits
\euro{}27/MWh on the heat pump's side, which shows that cavern access is not on its own
sufficient to give hydrogen the base-load case; Belgium and Ireland, which hold base
load without caverns, show that it is not strictly necessary either. Neither case
justifies a hydrogen heating programme, but both warrant the levy-incidence reform of
recommendation~3 before the base-load comparison in those countries is treated as
settled. Just Transition Fund resources should be linked explicitly to both grid
decarbonisation and building electrification, so that the two imperatives become
complementary.

\textbf{5. Policy should recognise that Switzerland's case for hydrogen-import
infrastructure for residential buildings is weak, and weak on delivered cost as well.}
The delivered-cost margin set out above is narrow, so it should not carry the argument
alone, and the grounds that bear on an irreversible national commitment all point the
same way.

Equalising the fiscal treatment of taxed Swiss retail electricity and an untaxed
imported molecule would widen the heat pump's margin. Hydrogen wins no Swiss winter
peak in any of the four scenarios on the reservoir-credited arm we read the Swiss
verdict from, on the figures given
above, so an import programme would not answer the winter adequacy problem that
electrified heat creates; and the same imported molecules
have higher-value uses in industry and mobility, where substitutes are fewer. Federal hydrogen strategy is therefore better directed towards industrial and mobility
applications, as the current strategy already does \citep{SFOE2023hydrogen}, while
residential resources are better directed towards heat-pump installation programmes,
grid flexibility, and seasonal thermal storage. This is offered as a direction following
from the base-load, peaking and capital-recovery results together and from competing
uses, and not as a modelled policy appraisal. The analysis also leaves out modelled
Swiss policy options, local infeasibility factors and other constraints.

%── 7.4 Future Research ─────────────────────────────────────────────────────────
\subsection{Directions for Future Research}

The study points to four productive directions for future work.

First, the least-cost optimisation layer currently fixes demand at the Stated
Policies trajectory and optimises the supply mix; a joint demand-and-supply
optimisation that co-optimises renovation depth with technology choice would close
the remaining gap between the simulation and optimisation frameworks. On the
supply side, the explicit unit-commitment dispatch introduced here
(Appendix~\ref{sec:results:dispatch}) now commits the firm fleet against ramp,
minimum-stable-output, reserve and start-up constraints, but it still relies on
calibrated availability profiles and a causal storage heuristic, not on
co-optimised storage and demand response. Endogenising storage sizing and
demand-side flexibility within the commitment model, and propagating the
capacity stack through the four policy scenarios, not a single central
trajectory, would further tighten the peaking and missing-money bounds.

Second, city-level spatial analysis within NUTS3 regions would capture
intra-regional heterogeneity that matters for the hydrogen case. Dense urban
cores with historic multi-family stock have feasibility profiles quite different
from those of peri-urban detached housing, and the aggregate NUTS3 treatment
masks this variation.

Third, the unit-commitment hourly dispatch quantifies one side of the
sector-coupling problem, namely the firm peaking capacity that electrified heat
requires for winter adequacy ($\approx$278~GW across the EU-27+UK+CH centrally,
about 31~per cent of the study-area system peak, with a wider band across the
interconnection and Dunkelflaute-depth sweep) and the
inability of energy-market revenue to finance it. A fully co-optimised treatment,
in which the increase in EU electricity demand from building electrification by
2050 is balanced against endogenous transmission, storage and flexibility
investment, would convert this adequacy signal into explicit grid-investment and
capacity-market requirements. This is also the natural setting in which to take up
the wider part of RQ3, the cross-sector synergies between buildings and industry
or transport that the present paper has deferred. A hydrogen turbine that cannot recover its capital
from heating-peak rent alone may yet be financeable when its capacity value is
shared across sectors, and quantifying that shared value is the central open
question the merit-order verdict leaves behind.

Fourth, within RQ2 a dedicated city-level analysis of Zurich, Geneva and Basel
would provide the granularity to distinguish hydrogen niches in the historic, dense
urban stock that the NUTS3 model treats as an aggregate, sharpening the RQ2 answer
for exactly the hard-to-treat districts where any residual case for hydrogen would
have to be made. This is a refinement of the existing scope rather than an extension of it.
